\documentclass[11pt]{article}
\usepackage[margin=1.1in]{geometry}
\usepackage{amsmath,amssymb,amsthm}
\usepackage{booktabs}
\usepackage[T1]{fontenc}
\usepackage{lmodern}
\usepackage{microtype}
\usepackage[hidelinks]{hyperref}
\usepackage{url}

\newtheorem{proposition}{Proposition}
\newtheorem{definition}[proposition]{Definition}
\newcommand{\HH}{\ensuremath{\mathcal{H}}}

\title{\bf The Information Order of Experiments:\\
Intervention Lattices and Generator Identifiability}
\author{
  Gregory J. Ward\thanks{Corresponding author: \texttt{greg@smartledger.technology}}
  \and Bryan W. Daugherty
  \and Shawn M. Ryan
}
\date{September 2026}

\begin{document}
\maketitle

\begin{abstract}
Which experiment should you run next? We give the question a finite, exactly computable
form. For a hypothesis space $\HH$ of candidate generators, write $O_1 \preceq O_2$ when
$O_1$ is obtainable from $O_2$ by garbling. Refinement then contracts preimages,
$P_{O_2}(g) \subseteq P_{O_1}(g)$, so per-generator ambiguity $i_O(g)$ and its population
mean $H(G \mid O)$ both fall---a one-line consequence, but one that turns ``more
informative'' into an order rather than a scale.

We instantiate the order exactly on all $1{,}500{,}625$ four-state partially observed
Markov chains, enumerating every subset of hidden-state resets. The reset sets form a
Boolean lattice under inclusion; informativeness itself is only a preorder, and we reserve
``lattice'' for the former. Four findings. The
passive experiment is a \emph{deterministic} coarsening of the full interventional one,
confirmed exactly, so its $1.874$ bits of shortfall is precisely
$I(G; O_I \mid O_P)$. Informativeness is genuinely partial, not total: a single reset
leaves $4.504$ bits where passive observation leaves $3.245$, so one intervention is
\emph{worse} than watching, while two are better. Which two matters enormously---two resets inside one observational
lump identify $0\%$ of the class, two that cross it identify $58.55\%$, at identical
cost---but crossing is not the rule. A sweep over all $15$ partitions refutes it, and the
exact criterion is conditional information: the value of adding reset $j$ to a set $A$ is
$H(Z_j \mid Z_A) = H(Z_j) - I(Z_j; Z_A)$, intrinsic information less what is already
known. And information gain is
submodular across every triple we can form, which places greedy experiment selection on
familiar footing.

The organising claim is that resolving a mechanism is partition refinement, that the order
on experiments controls its rate, and that the value of an intervention is the conditional
information it contributes beyond the experiments already performed.
\end{abstract}

\section{Introduction}

A companion paper~\cite{paper1} argues that identifiability should be measured rather than
assumed, by counting the preimage $P_O(g) = \{h \in \HH : O(h) = O(g)\}$ of an observation
map, and reports that identifiability and predictive uncertainty are distinct coordinates.
It leaves a question open. Identifiability there is relative to a declared observation map
and a declared class of admissible interventions, and changing that class moved $82\%$ of
one hypothesis space between regimes. If knowability is regime-relative, the regimes
themselves need a structure.

They have one. Experiments are partially ordered by informativeness, and the order is old:
Blackwell's comparison of experiments~\cite{blackwell1953} ranks one experiment above
another when the second can be obtained from the first by a stochastic garbling. What we
add is a finite, exactly enumerable instance in which the whole lattice can be computed
rather than argued about, together with three empirical facts that are not consequences of
the order alone.

\section{The refinement order}

Let $\HH$ be finite, $G$ uniform on $\HH$, and let $O_1, O_2$ be observation maps.

\begin{definition}
$O_1 \preceq O_2$ iff there is a map $F$ with $O_1 = F \circ O_2$.
\end{definition}

\begin{proposition}[Containment]\label{prop:contain}
If $O_1 \preceq O_2$ then for every $g \in \HH$,
$P_{O_2}(g) \subseteq P_{O_1}(g)$, hence $i_{O_2}(g) \le i_{O_1}(g)$ and
$H(G \mid O_2) \le H(G \mid O_1)$.
\end{proposition}
\begin{proof}
If $O_2(h) = O_2(g)$ then $O_1(h) = F(O_2(h)) = F(O_2(g)) = O_1(g)$. Containment gives the
inequality on $\log_2$ of cardinalities pointwise, and averaging preserves it.
\end{proof}

\begin{proposition}[Value of refinement]\label{prop:value}
If $O_1 \preceq O_2$ then
$H(G \mid O_1) - H(G \mid O_2) = I(G; O_2 \mid O_1)$.
\end{proposition}

Proposition~\ref{prop:value} is what licenses reading a difference of conditional
entropies as the information the finer experiment carries that the coarser one cannot
reach. It requires the order; between incomparable experiments the difference is a
comparison of two numbers and nothing more.

\begin{proposition}[Submodularity]\label{prop:submod}
Let $Z_j = O_j(G)$ for $j$ in a finite index set $J$, with each $Z_j$ a deterministic
function of $G$, and write $Z_A = (Z_j)_{j \in A}$. Then
$F(A) = I(G; Z_A) = H(Z_A)$ is monotone and submodular in $A$. The same holds
conditionally on any fixed observation $P$: $F_P(A) = I(G; Z_A \mid P) = H(Z_A \mid P)$
is monotone and submodular.
\end{proposition}
\begin{proof}
$H(Z_A \mid G) = 0$ because each $Z_j$ is a function of $G$, so
$I(G; Z_A) = H(Z_A) - H(Z_A \mid G) = H(Z_A)$. The entropy of a collection of random
variables is a monotone submodular set function, which is Shannon's inequality
$H(Z_{A \cup \{j\}}) - H(Z_A) \ge H(Z_{B \cup \{j\}}) - H(Z_B)$ for $A \subseteq B$.
Conditioning on a fixed $P$ preserves both properties, and
$I(G; Z_A \mid P) = H(Z_A \mid P)$ by the same argument.
\end{proof}

The pair $(J, r)$ with $r(A) = H(Z_A)$ is therefore an \emph{entropic polymatroid}, the
dependence structure Fujishige identified for collections of random
variables~\cite{fujishige1978}. This is more specific than ``some submodular objective'':
$r(\{j\}) = H(Z_j)$ is the standalone information of experiment $j$,
$r(A \cup \{j\}) - r(A) = H(Z_j \mid Z_A)$ its marginal information, and
$r(\{i\}) + r(\{j\}) - r(\{i,j\}) = I(Z_i; Z_j)$ the pairwise redundancy. Our contribution
is not the greedy machinery, which is standard in observation
selection~\cite{krause2005}, but that the finite-generator intervention problem induces an
explicit entropic polymatroid whose rank values are exactly computable from generator
equivalence classes, making those guarantees directly applicable.

\begin{proposition}[One-way migration]\label{prop:oneway}
Classify generators by thresholds on $i_O(g)$ and on a behavioural coordinate $p_k(g)$.
If $p_k$ is evaluated under an operating condition held fixed across $O_1 \preceq O_2$,
then under refinement no generator can become less identifiable, and none can move in the
behavioural coordinate.
\end{proposition}

This is worth stating because it bounds what a migration statistic can mean. Reported
migration under refinement is necessarily one-dimensional and one-directional; it is not
evidence of two-dimensional mobility.

\section{The benchmark system}

We reuse the family of~\cite[\S9]{paper1} rather than introduce another: hidden state
$X_t \in \{0,1,2,3\}$, transition rows drawn from
$\{0,\frac14,\frac12,\frac34,1\}$ summing to one ($35$ rows), all $35^4 = 1{,}500{,}625$
generators enumerated completely, observation lumping $O(0)=O(1)=0$, $O(2)=O(3)=1$,
horizon $T=6$. The passive experiment starts from the uniform hidden state. For
$A \subseteq \{0,1,2,3\}$, the experiment $O_A$ resets the hidden state to each $j \in A$
in turn and returns the resulting observable laws.

\section{Results}

\paragraph{The passive experiment is a deterministic coarsening.}
Writing $L_j$ for the law obtained by resetting to state $j$, we verified exactly, over
every generator, that
\[
  O_P \;=\; \tfrac14 \textstyle\sum_{j} L_j ,
\]
so $O_P \preceq O_{\{0,1,2,3\}}$ with a deterministic $F$. Proposition~\ref{prop:contain}
then applies, and we confirm its conclusion computationally as an implementation check:
zero containment violations across every pair $A \subset B$, and zero against the passive
mixture.

\begin{table}[h]
\centering
\begin{tabular}{lrrrr}
\toprule
Reset set $A$ & $H(G \mid O_A)$ & classes & $\Pr[i_{O_A}=0]$ & $\max |P|$ \\
\midrule
$\varnothing$ & 20.517 & 1 & $0.00\%$ & 1{,}500{,}625 \\
any singleton & 4.504 & 400{,}584 & $0.00\%$ & 67{,}375 \\
$\{0,1\}$ or $\{2,3\}$ & 3.356 & 482{,}902 & $0.00\%$ & 30{,}625 \\
$\{0,2\},\{0,3\},\{1,2\},\{1,3\}$ & 2.618 & 967{,}107 & $58.55\%$ & 8{,}953 \\
any triple & 1.914 & 1{,}101{,}899 & $69.84\%$ & 7{,}263 \\
$\{0,1,2,3\}$ & 1.371 & 1{,}235{,}809 & $82.01\%$ & 6{,}561 \\
\midrule
passive (mixture) & 3.245 & 303{,}229 & $0.00\%$ & 7{,}537 \\
\bottomrule
\end{tabular}
\caption{The intervention lattice, exact over $1{,}500{,}625$ generators. Rows within a
group are identical to the digits shown, reflecting the symmetry of the construction.}
\label{tab:lattice}
\end{table}

\paragraph{Informativeness is partial, not total.}
A single reset leaves $4.504$ bits; passive observation leaves $3.245$. One intervention
is \emph{less} identifying than watching a system you did not touch. There is no
contradiction with Proposition~\ref{prop:contain}: $O_P$ is a coarsening of the full reset
set, not of any singleton, so the two are incomparable in $\preceq$ and the proposition
says nothing about them. Two resets ($2.618$) overtake passive; one does not. Three relations are in play and
only two are comparabilities:
\[
  O_P \preceq O_{\{0,1,2,3\}}, \qquad
  O_{\{0\}} \preceq O_{\{0,1,2,3\}}, \qquad
  O_{\{0\}} \;\;\text{and}\;\; O_P \;\;\text{incomparable.}
\]
Passive is a garbling of the full reset protocol, but neither passive nor a single reset is
a garbling of the other, so their entropies may fall in either order and here they fall the
surprising way. Any claim of the form ``intervention beats observation'' is false as
stated; the correct statement is that informativeness is a partial order.

\paragraph{The interventions that pay are those crossing the observation's kernel.}
Let $x \sim_O y$ iff $O(x) = O(y)$, here $\{0,1\}$ and $\{2,3\}$. Two resets chosen inside
one $\sim_O$ class leave $3.356$ bits and identify \emph{nothing}; two that cross leave
$2.618$ bits and uniquely identify $58.55\%$ of the entire hypothesis class. The cost is
identical. Greedy selection, run without knowledge of the lumping, chooses $0$, then $2$,
then $3$, then $1$---crossing at the first opportunity.

We stated this as a conjecture---interventions crossing $\sim_O$ dominate interventions of
equal cardinality confined within one class---and then tested it, with criteria fixed in
advance, over \emph{all} $15$ set partitions of the four hidden states. It is false.

The sweep uses a coarser but still completely enumerated family so that every partition is
affordable: transition rows from $\{0,\frac12,1\}$, giving $10^4 = 10{,}000$ generators, at
$T=5$; entropies are therefore not comparable in absolute terms with
Table~\ref{tab:lattice}. Within each (partition, $|A|$) cell we asked whether mean
$H(G \mid O_A)$ is non-increasing in the number of blocks $A$ meets. It is in only
$3$ of $23$ cells---$13\%$, against a pre-committed threshold of $90\%$.

The failures have a clean shape. For the partition $\{0\} \mid \{1,2,3\}$ the effect
\emph{reverses}: two resets confined to the large block leave $4.852$ bits, while two that
cross leave $5.917$---crossing is worse by $1.065$ bits. For the balanced
$\{0,1\} \mid \{2,3\}$ the original direction holds, $3.799$ against $4.187$. Restricting
the test to partitions whose blocks are all the same size, it passes $3/3$.

No replacement conjecture is required, because the exact rule is already available.
Since each $Z_j$ is a deterministic function of $G$, Proposition~\ref{prop:submod} gives
$F(A) = H(Z_A)$, so the value of adding reset $j$ to a set $A$ is exactly its conditional
novelty
\[
  \Delta(j \mid A) \;=\; H(Z_j \mid Z_A),
\]
and for a pair, $H(Z_i,Z_j) = H(Z_i) + H(Z_j) - I(Z_i;Z_j)$. A pair of interventions is
valuable when both carry information and little of it is shared. Block crossing was only
ever a proxy for low redundancy, and decomposing the two cases (Experiment~25) shows they
are governed by different terms.

\begin{table}[h]
\centering
\begin{tabular}{llrrrr}
\toprule
Partition & pair & crossing & $H(Z_i){+}H(Z_j)$ & $I(Z_i;Z_j)$ & $H(G \mid O_A)$ \\
\midrule
$\{0\}\mid\{1,2,3\}$ & $\{1,2\}$ & no  & 12.056 & 3.621 & 4.852 \\
                       & $\{0,1\}$ & yes & 11.045 & 3.675 & 5.917 \\
\addlinespace
$\{0,1\}\mid\{2,3\}$ & $\{0,1\}$ & no  & 14.342 & 5.241 & 4.187 \\
                        & $\{0,2\}$ & yes & 14.342 & 4.854 & 3.799 \\
\bottomrule
\end{tabular}
\caption{The crossing effect and its reversal, decomposed by
$H(Z_i,Z_j) = H(Z_i) + H(Z_j) - I(Z_i;Z_j)$.}
\label{tab:decomp}
\end{table}

On the balanced partition the individual entropies are \emph{exactly equal}, so the entire
$0.387$-bit advantage of crossing is a redundancy effect: crossing pairs share less
($4.854$ against $5.241$). On the unbalanced partition the redundancies are nearly equal
---and the crossing pair is marginally the \emph{more} redundant of the two---so
essentially all of the $1.065$-bit gap, $1.011$ of it, comes from the individual term: a
reset into a singleton observational block carries only $5.017$ bits where a reset into the
three-state block carries $6.028$.

We had guessed the singleton reset was redundant. It is not; it is simply less informative
on its own, which is a different mechanism and one the decomposition settles rather than
argues. Crossing is a heuristic for low redundancy that works when symmetry equalises the
individual terms and fails when it does not.

The selection rule that follows is exact rather than heuristic:
\[
  j^\star \;=\; \operatorname*{arg\,max}_{j \notin A}\; H(Z_j \mid Z_A),
\]
which is the greedy step of Proposition~\ref{prop:submod} written in the variables that
make it computable. What remains genuinely open is not the rule but its explanation:
\emph{which structural features of $(\HH, O)$ make $H(Z_j \mid Z_A)$ large?} Observation
block size re-enters there as a candidate explanatory variable rather than as the theory.
A natural refinement to study is the interventional signature partition
$x \sim_Z y \iff Z_x = Z_y$, which is not the observational partition $\sim_O$: states
alike under $\sim_O$ but distinct under $\sim_Z$ are exactly the distinctions intervention
can unlock, and states alike under both are hidden from this experimental family
altogether.

We record the refutation rather than quietly narrowing the claim, because the conjecture as
first written was the natural reading of Table~\ref{tab:lattice} and would have been an
easy thing to leave standing.

\paragraph{Information gain is submodular---by Proposition~\ref{prop:submod}, not by
observation.}
An earlier version of this paper reported $108$ submodularity inequalities holding with
zero violations and described the property as measured. That was a category error. In this
setting the observable law of each reset is an exact deterministic function of the
generator, so $F(A) = H(Z_A)$ and submodularity is Shannon's inequality; the computation
could not have failed, and it checks our implementation rather than the mathematics. We
record it as such: $108/108$, consistent with the theorem.

What the theorem buys is operational. Maximising a monotone submodular function under a
cardinality constraint admits the classical $(1 - 1/e)$ greedy
guarantee~\cite{nemhauser1978}, so under a budget $|A| \le k$ the rule
\[
  j^\star_t \;=\; \operatorname*{arg\,max}_{j \notin A_t}\; I\big(G; Z_j \mid P, Z_{A_t}\big)
\]
---perform next the experiment expected to eliminate the most remaining ambiguity---is
within $1 - 1/e$ of the best possible set of that size. That holds for any family of
experiments whose exact outputs are deterministic in the unknown generator, which is the
regime this framework assumes throughout. The greedy gain sequence here is $16.013$,
$1.886$, $0.704$, $0.543$ bits; the fourth reset is nearly redundant.

\paragraph{The polymatroid has no non-trivial closure.}
A polymatroid carries a closure operator, $j \in \operatorname{cl}(A)$ iff
$H(Z_j \mid Z_A) = 0$, and a basis is a minimal spanning set. We computed these exactly
---$H(Z_J \mid Z_A) = 0$ holds iff $A$ induces as many generator classes as $J$, an
integer comparison rather than a tolerance---for all $14$ non-trivial partitions of the
four hidden states. In every one, the \emph{only} spanning set is $J$ itself. There is no
experimental compression here: no reset is redundant given the other three, and a
four-reset protocol cannot be replaced by a smaller one.

The margins by which this fails are informative. $H(Z_j \mid Z_{J \setminus \{j\}})$ ranges
from $0.280$ to $0.902$ bits, and within every partition it is monotone in the size of
$j$'s observation block, with zero violations across all $14$: a reset into a state the
observation already isolates contributes least that the others cannot supply. Notably
block size does \emph{not} predict $H(Z_j)$ itself---Experiment~25 found those means
non-monotone at $8.050$, $8.157$, $6.028$ for block sizes $1$, $2$, $3$---so the structural
variable governs the conditional contribution and not the standalone one. This is the
defensible residue of the intuition that failed in Experiment~24, and it is worth noting
that the pairwise redundancies did not reveal it: $I(Z_0;Z_1)$ was marginally
\emph{lower} for the singleton in the reversal case, while the conditional-on-all-others
measure separates cleanly.

A further contrast: the rank values $r(J)$ vary widely with the observation partition,
from $9.892$ to $13.288$ bits, while the basis structure does not vary at all. What the
observation map changes is how much the experiments tell you, not which of them are
needed.

\paragraph{Compression exists, but it is a property of the hypothesis class.}
The absence of compression above is suspicious, because $\HH = \mathrm{ROWS}^4$ is a
product: the four transition rows are chosen independently, which is exactly the structure
that should resist it. We therefore compared three classes on the same system and the same
$14$ partitions, with criteria fixed in advance. Alongside the product class we took a
\emph{tied} class constraining row $3$ to equal row $0$, giving $10^3$ generators, and a
\emph{random} subset of the product class of the same cardinality. The random class is the
control: it decides whether any compression is structural or merely an effect of having
fewer generators to tell apart.

\begin{table}[h]
\centering
\begin{tabular}{lrrl}
\toprule
Hypothesis class & $|\HH|$ & partitions compressing & minimal bases \\
\midrule
product, rows free            & 10{,}000 & $0/14$  & $\{0,1,2,3\}$ \\
tied, row $3 =$ row $0$       & 1{,}000  & $14/14$ & $\{0,1,2\}$, $\{1,2,3\}$ \\
random subset, same size      & 1{,}000  & $0/14$  & $\{0,1,2,3\}$ \\
\bottomrule
\end{tabular}
\caption{Compression is structural, not a cardinality effect: the tied class compresses in
every partition while a random class of identical size compresses in none.}
\label{tab:structure}
\end{table}

The tied class compresses in every partition, always to a basis of size three, and always
by omitting one of the two tied states---never states $1$ or $2$, whose rows remain free.
The equally sized random class compresses nowhere. So Experiment~26's result is a statement
about the hypothesis class, not about interventions: \emph{an intervention is redundant
exactly when the generators carry a dependency that makes it so.}

It is tempting to read this as a clean separation---observation sets the rank magnitudes,
hypothesis-class structure sets the basis---and an earlier version of this paper said so.
Experiment~28 refutes it. Repeating the comparison across six classes and all $14$
partitions, the basis is invariant under the observation map for the equality and adjacent
ties, but \emph{varies} for a permutation tie $R_3 = \pi(R_0)$ and for a two-row derived
constraint $R_3 = f(R_0, R_1)$. So the observation map can change basis structure, for some
dependency types and not others, and characterising which is open.

\paragraph{Which generator dependencies survive the observation map.}
A constraint among generator coordinates need not become a functional dependency among
reset outputs, because $Z_j$ is a function of the whole generator and not of row $j$ alone.
We made this precise with two closure operators, both tested by exact integer class counts:
$\operatorname{cl}_G(A) = \{j : H(R_j \mid R_A) = 0\}$ among generator coordinates and
$\operatorname{cl}_Z(A) = \{j : H(Z_j \mid Z_A) = 0\}$ among reset outputs.

\begin{table}[h]
\centering
\begin{tabular}{llrll}
\toprule
Class & constraint & $|\HH|$ & minimal bases & basis invariant in $O$ \\
\midrule
FREE     & none                        & 10{,}000 & $\{0,1,2,3\}$ & yes \\
EQUAL    & $R_3 = R_0$                 & 1{,}000  & $\{0,1,2\},\{1,2,3\}$ & yes \\
PERM     & $R_3 = \pi(R_0)$            & 1{,}000  & $\{0,1,2\},\{1,2,3\}$ & \textbf{no} \\
DERIVED  & $R_3 = f(R_0,R_1)$          & 1{,}000  & $\{0,1,2\},\{0,2,3\},\{1,2,3\}$ & \textbf{no} \\
ADJACENT & $R_2 = R_1$                 & 1{,}000  & $\{0,1,3\},\{0,2,3\}$ & yes \\
RANDOM   & size-matched control        & 1{,}000  & $\{0,1,2,3\}$ & yes \\
\bottomrule
\end{tabular}
\caption{Six hypothesis classes, five at identical cardinality. Which rows are tied
determines which resets become droppable; a two-row constraint yields three bases rather
than two.}
\label{tab:dependency}
\end{table}

Two pre-registered criteria failed, and both failures are informative. The inclusion
$\operatorname{cl}_G(A) \subseteq \operatorname{cl}_Z(A)$ is \emph{false}: $181$
violations. The mechanism is transparent once seen---under $R_3 = R_0$ we have
$3 \in \operatorname{cl}_G(\{0\})$, but $Z_3$ also depends on rows $1$ and $2$, so $Z_0$
alone cannot determine it. A coordinate dependency is not an output dependency.

The converse count was $0$ in these six classes, and we briefly took
$\operatorname{cl}_Z \subseteq \operatorname{cl}_G$ for an empirical regularity. It is
not one. For deterministic experiment maps in general it fails at once: with $R_0, R_1$
independent and $Z_0 = Z_1 = R_0 \oplus R_1$, we have $1 \in \operatorname{cl}_Z(\{0\})$
while $1 \notin \operatorname{cl}_G(\{0\})$. Experiment~29 constructs the same effect
inside the reset-law architecture.

\paragraph{The observation map can create dependencies: lumpability.}
Fix the balanced partition and constrain each row so its \emph{block} transition sums
depend only on the source state's block---states $0$ and $1$ both send mass $a$ to block
$0$, states $2$ and $3$ both send $b$---while the split within each block stays free. Then
$R_0$ and $R_1$ are not functions of one another, yet the observed process is Markov on
blocks and resets $0$ and $1$ induce the same observable law. Characterising each class by
$(|\mathcal{P}|, |\mathcal{D}|, |\mathcal{C}|)$---dependencies preserved, destroyed and
created---under that partition:

\begin{table}[h]
\centering
\begin{tabular}{llrrrr}
\toprule
Class & constraint & $|\HH|$ & $|\mathcal{P}|$ & $|\mathcal{D}|$ & $|\mathcal{C}|$ \\
\midrule
FREE          & none                        & 10{,}000 & 0 & 0 & 0 \\
TIED          & $R_3 = R_0$                 & 1{,}000  & 8 & 0 & 0 \\
LUMPABLE      & block sums tied by block    & 1{,}156  & 0 & 0 & \textbf{16} \\
TIED\_LUMPABLE & both                       & 340      & 8 & 0 & \textbf{8} \\
RANDOM        & size-matched control        & 1{,}156  & 0 & 0 & 0 \\
\bottomrule
\end{tabular}
\caption{Redundancy has at least two independent origins. The lumpable class creates $16$
dependencies with no generator dependency to inherit; a random class of the same
cardinality creates none.}
\label{tab:pdc}
\end{table}

The created dependencies are exactly $Z_0 \leftrightarrow Z_1$ and
$Z_2 \leftrightarrow Z_3$ and their extensions, and they appear under \emph{one} of the
$14$ partitions---the balanced one the class was built for. Under any other observation the
lumpability does not hold and nothing is created. The effect is the quotient matching the
constraint, not the constraint alone.

\paragraph{A criterion for survival.}
The robust/fragile split is a taxonomy, not a criterion. Decomposing a reset law by its
first symbol,
\[
  Z_s^{(T)} \;=\; \delta_{O(s)} \otimes \sum_x P_{sx}\, Z_x^{(T-1)},
\]
makes the mechanism visible and suggests one. Under an identity tie $R_i = R_j$ the two
sums agree term by term for every $O$ and every $T$, so robustness is provable rather than
observed; we ran it as an implementation check and found $0$ counterexamples over $2002$
generator--partition pairs. Under a permutation tie $R_j = \pi(R_i)$ the comparison is
$\sum_y P_{iy} Z_{\pi(y)}^{(T-1)}$ against $\sum_y P_{iy} Z_y^{(T-1)}$, which agree when
$\pi$ moves each \emph{reachable} destination to a behaviourally equivalent state, where
$s \sim t$ iff $Z_s^{(T-1)} = Z_t^{(T-1)}$ for that generator. Tested in the sufficient
direction over $2002$ pairs:

\begin{center}
\begin{tabular}{lrr}
\toprule
 & suffix laws equal & suffix laws differ \\
\midrule
$\pi$ behaviourally compatible   & 153 & \textbf{0} \\
$\pi$ incompatible               & 358 & 1491 \\
\bottomrule
\end{tabular}
\end{center}

Zero counterexamples: compatibility suffices. It is not necessary---$358$ cases agree
without it---which is expected, since the sums can coincide numerically without the
permutation respecting behaviour.

For the created case we pre-registered that dependencies appear exactly on partitions
satisfying strong lumpability. That \emph{failed}, on one partition: the discrete
partition $0|1|2|3$ is strongly lumpable only vacuously, every block being a singleton, and
creates nothing because there is nothing for the quotient to collapse. The refinement
---lumpable \emph{and} some block of size greater than one---matches dependency creation
on all $14$ partitions, and we record it as post-hoc.

\paragraph{One equation underneath all three.}
The three mechanisms are three cases of a single linear condition. Let $b_x$ be the
observable law of the next $T-1$ symbols from hidden state $x$, and $B_{O,T-1}(G)$ the
matrix with those rows. Resetting to $s$ and taking one transition leaves the hidden
distribution $P_s$, so the suffix law is $S_s^{(T)} = P_s B$.

\begin{proposition}[Semantic kernel]\label{prop:kernel}
For a finite generator $G$, observation map $O$ and horizon $T$,
\[
  S_s^{(T)} = S_t^{(T)}
  \quad\Longleftrightarrow\quad
  (P_s - P_t)\,B_{O,T-1}(G) = 0
  \quad\Longleftrightarrow\quad
  P_s - P_t \in \ker B_{O,T-1}(G).
\]
\end{proposition}
\begin{proof}
$S_s = P_s B$ and $S_t = P_t B$, so $S_s = S_t$ iff $(P_s - P_t)B = 0$.
\end{proof}

Each earlier mechanism is a case. An \emph{identity tie} has $P_s - P_t = 0$, which lies
in every kernel---so its robustness across all $O$ and $T$ is immediate, and needs no
appeal to behaviour. A \emph{permutation tie} $P_t = P_s \Pi$ survives iff
$P_s(I - \Pi)B = 0$, which is strictly weaker than the behavioural compatibility
$\Pi B = B$ we tested in Experiment~30: compatibility forces each moved state's law to
agree separately, whereas the kernel condition needs only their $P_s$-weighted mixture to
cancel. That is what the $358$ agreements without compatibility were---weighted
cancellation, not coincidence. \emph{Lumpability} is the case $P_i \ne P_j$ with
$P_i - P_j \in \ker B$: the rows differ in a direction the observation is blind to.
Verified over $11{,}004$ reset pairs with zero failures, as an implementation check.

\paragraph{Blindness is structural or transient.}
$\dim \ker B$ measures how many directions of hidden transition behaviour the observation
cannot see, and it separates the classes sharply. Under the discrete partition, where every
state is its own block, it is $0$ for every generator of every class: full observation is
blind to nothing. Under the balanced partition the lumpable class has $\dim \ker B = 2$
for all $32$ generators tested, at horizons $T = 3$, $5$ and $7$ alike, and the kernel is
spanned by exactly $e_0 - e_1$ and $e_2 - e_3$---the within-block differences, one per
block beyond the first member.

The contrast with an unconstrained class is the point. For FREE generators under the same
partition the kernel dimensions are distributed $\{2\!:\!6,\ 1\!:\!17,\ 0\!:\!10\}$ at
$T=3$ and contract to $\{2\!:\!6,\ 1\!:\!14,\ 0\!:\!13\}$ by $T=5$, then stop. So
kernels contract with horizon---non-increasing across $T = 2, \dots, 6$ with zero
exceptions---but the lumpable kernel does not contract at all. There are two kinds of
observational blindness: \emph{transient}, removable by watching longer until it settles,
and \emph{structural}, which no horizon removes. That is the linear-algebraic counterpart
of the preimage stabilisation at $T^\star$ reported in~\cite[\S9]{paper1}: preimages
contract over candidate generators, kernels contract over hidden-state directions, and both
reach a floor.

Proposition~\ref{prop:kernel} also gives the robust/fragile distinction a definition rather
than a label: a dependency is \emph{robust} when its difference vector lies in
$\ker B_{O,T}$ for every admissible $O$, which identity ties satisfy because that vector is
zero, and \emph{fragile} when membership varies with $O$.

\paragraph{Robust and fragile dependencies.}
The two failures of Experiment~28 turn out to be one phenomenon. Breaking its $181$
inclusion violations down by class gives $100$ for PERM and $81$ for DERIVED, and zero for
every other class---and those are precisely the two classes whose bases varied with the
observation map. A dependency that survives every observation map cannot move the basis; a
dependency some observation maps destroy must. So the constraint types split into
\emph{robust} (the identity ties EQUAL and ADJACENT: nothing destroyed, basis invariant)
and \emph{fragile} (PERM and DERIVED: dependencies destroyed, basis varies), with
lumpability supplying a third category that owes nothing to the generator at all.

\paragraph{The value of full intervention.}
$H(G \mid O_P) - H(G \mid O_I) = 3.245 - 1.371 = 1.874$ bits, which by
Proposition~\ref{prop:value} is exactly $I(G; O_I \mid O_P)$: the generator information the
complete reset protocol carries that no amount of passive observation can reach.

\section{Related work}

The order is Blackwell's~\cite{blackwell1953}; our instance is its degenerate deterministic
case. Choosing experiments to maximise expected information is Lindley's~\cite{lindley1956},
and choosing interventions to secure identifiability is well developed in causal discovery,
where optimal strategies are known~\cite{hauser2014}. The greedy guarantee for submodular
objectives is~\cite{nemhauser1978}. What we contribute is not a new criterion but an
exactly enumerated lattice: every experiment in a finite family, every containment relation
checked rather than assumed, and the resulting empirical facts---incomparability, the
crossing rule, submodularity---measured rather than argued.

\section{Limitations}

The lattice is computed on one family with one lumping, and the striking symmetry of
Table~\ref{tab:lattice} is a property of that construction. We do not know what distinguishes robust from fragile dependencies beyond the four
examples in Table~\ref{tab:dependency}; identity ties are robust and permutation or
two-row constraints are fragile, which is a pattern with $n = 4$. The lumpable
construction of Table~\ref{tab:pdc} is one aggregation matched to one partition, and we
have not characterised which quotients create dependencies for which classes. The crossing conjecture is refuted in general and survives only for equal-block
partitions; the exact selection rule $\Delta(j \mid A) = H(Z_j \mid Z_A)$ replaces it, but we have no
structural characterisation of when that quantity is large, which is the open problem this
paper leaves. Submodularity is a theorem here (Proposition~\ref{prop:submod}) and holds whenever exact
experiment outputs are deterministic functions of the generator; it would need revisiting
for noisy or sampled observations, where $H(Z_A \mid G) > 0$ and the identity
$F(A) = H(Z_A)$ fails. Identification is at horizon $T=6$, which is where the partition of this family
stabilises~\cite[\S9]{paper1}, but stability over a finite window is evidence rather than
proof. Proposition~\ref{prop:contain} is a theorem, so the computation confirming it tests
our implementation, not the mathematics.

\section{Discussion}

Resolving a mechanism is partition refinement. Beginning with $[g] = \HH$, each experiment
induces an equivalence relation and the cell containing the truth contracts,
\[
  \HH \supseteq [g]_{\Pi_1} \supseteq [g]_{\Pi_2} \supseteq \cdots ,
\]
with perfect resolution at $[g]_\Pi = \{g\}$ and irreducible ambiguity when the admissible
order bottoms out while $|[g]_\Pi| > 1$. Three structures sit on top of one another:
generator space is partitioned by an observation; the information order says which
partitions refine which; and submodularity controls how efficiently refinements can be
bought under a budget. What the measurements add is that the rate of
contraction is controlled by the order on experiments, that the order is partial rather
than total---so more control is not uniformly more information---and that the experiments
worth buying are the ones that cross the distinctions the observation map erases.

\paragraph{Reproducibility.}
\texttt{mgs\_lattice.py} produces the intervention-order figures and
\texttt{mgs\_partitions.py} the partition sweep of \S4; \texttt{audit.py} re-derives every
number here against them, and reports its status in \texttt{AUDIT.txt}. All are in the
repository accompanying~\cite{paper1}.

\begin{thebibliography}{9}
\bibitem{blackwell1953} D. Blackwell. Equivalent comparisons of experiments.
\emph{Annals of Mathematical Statistics}, 24(2):265--272, 1953.
\bibitem{fujishige1978} S. Fujishige. Polymatroidal dependence structure of a set of
random variables. \emph{Information and Control}, 39(1):55--72, 1978.
\bibitem{hauser2014} A. Hauser and P. B\"uhlmann. Two optimal strategies for active
learning of causal models from interventional data. \emph{International Journal of
Approximate Reasoning}, 55(4):926--939, 2014.
\bibitem{krause2005} A. Krause and C. Guestrin. Near-optimal nonmyopic value of
information in graphical models. In \emph{Proceedings of the 21st Conference on
Uncertainty in Artificial Intelligence}, 2005.
\bibitem{lindley1956} D. V. Lindley. On a measure of the information provided by an
experiment. \emph{Annals of Mathematical Statistics}, 27(4):986--1005, 1956.
\bibitem{nemhauser1978} G. L. Nemhauser, L. A. Wolsey, and M. L. Fisher. An analysis of
approximations for maximizing submodular set functions---I. \emph{Mathematical
Programming}, 14(1):265--294, 1978.
\bibitem{paper1} G. J. Ward, B. W. Daugherty, and S. M. Ryan. Identifiability and
predictability are distinct: measured preimages across generative systems, 2026.
Unpublished draft.
\end{thebibliography}

\end{document}
